\chapter{Persistencia de desborde}
\label{ch:overflow}

\begin{note}
Stub. Cuando el pipeline de ingesta se cancela (el usuario cierra
la sesión, el proceso termina) o cuando el \code{parsed\_channel}
se llena más rápido de lo que el writer puede drenar, los
\code{ParsedBatch}es encolados se persisten a disco en vez de
descartarse. En la próxima carga de sesión el desborde se
reproduce en orden antes de que el streaming se reanude.
\end{note}

\section{Invariantes}

\begin{itemize}
    \item Cada batch persistido retiene su \code{source\_id},
        \code{watermark} y prioridad, así que la reanudación
        avanza el watermark de la fuente exactamente como lo
        haría una ingesta en caliente.
    \item El replay es idempotente: el grafo deduplica por
        $(source, dest, rel, timestamp)$ en
        \code{GraphWriter::insert\_triples\_chunk}, así que el
        mismo batch aplicado dos veces tiene el mismo efecto que
        una sola.
\end{itemize}

\section{Formato de archivo}

Los batches de desborde se guardan bajo el directorio de sesión
como archivos JSON. El formato es estable entre reinicios y
deliberadamente legible para humanos: un operador puede
inspeccionar el trabajo encolado con \code{jq} sin más herramientas
que un editor de texto.

\begin{inthecode}
\filepath{graph\_hunter\_core/src/ingest/overflow.rs} ---
persistencia y replay.
\end{inthecode}
